\section{Where the interventions fail}
\label{sec:failure}

A harness change that only ever helps would be suspicious. Three regimes where
ours do not, or where the phenomenon itself weakens, are worth stating plainly.

\paragraph{The ban displaces repetition into paraphrase.}
Blocking a token sequence guarantees the model cannot emit that sequence; it
does not make the model change its mind. In rollouts we observe the predicted
consequence directly: a model whose exact repeat is blocked will sometimes
produce a near-identical call with a reordered argument, different quoting, or a
trivially altered value. This is why we report the canonical repeat rate, which
counts those as repeats, alongside the exact rate. The gap between the two
metrics is the honest measure of how much the ban actually accomplishes, and it
is the reason we do not recommend the ban on its own.

Figure~\ref{fig:qualitative-ban} shows a pair of trajectories from the same
task and model under \emph{verbatim} and \emph{verbatim+ban}, and it is worth
reading closely because it makes the mechanism visible in a way the aggregates
cannot. Under the standard harness the model emits one wrong call and then
repeats it, byte for byte, five times. Under the ban it cannot, and what it
does instead is not to reconsider the call: it keeps the same call and starts
appending fragments of the error message to it as an extra argument
(\texttt{error='write\_file: SyntaxError: could not parse'}). The exact repeat
rate for that rollout is zero, and nothing has been fixed. Whatever is
producing this behaviour is copying from the transcript, and forbidding one
particular string from the transcript redirects it to the nearest other string
rather than turning it off.

\begin{figure*}[t]
\centering
\begin{minipage}{\textwidth}
\textbf{qwen2.5-0.5b, harness=verbatim}
\begin{scriptsize}\begin{verbatim}
goal: agent-0011  solved=False
1. [ok ] convert_units(value=88, from_unit='km', to_unit='mi')
2. [ERR] write_file(path='reports/conversion.md', content=str(54.6806))
3. [ERR] write_file(path='reports/conversion.md', content=str(54.6806))
4. [ERR] write_file(path='reports/conversion.md', content=str(54.6806))
5. [ERR] write_file(path='reports/conversion.md', content=str(54.6806))
6. [ERR] write_file(path='reports/conversion.md', content=str(54.6806))
\end{verbatim}\end{scriptsize}

\textbf{qwen2.5-0.5b, harness=verbatim-ban}
\begin{scriptsize}\begin{verbatim}
goal: agent-0011  solved=False
1. [ok ] convert_units(value=88, from_unit='km', to_unit='mi')
2. [ERR] write_file(path='reports/conversion.md', content=str(54.6806))
3. [ERR] write_file(path='reports/conversion.md', content=str(54.6806), error='write_file: argument value ~
4. [ERR] write_file(path='reports/conversion.md', content=str(54.6806), error='write_file: SyntaxError: co~
5. [ERR] write_file(path='reports/conversion.md', content=str(54.6806), error='write_file: SyntaxError: co~
6. [ERR] write_file(path='reports/conversion.md', content=str(54.6806), error='write_file: argument value ~
\end{verbatim}\end{scriptsize}
\end{minipage}
\caption{The same task and model under the standard harness and under the decoder ban. Selected by rule, not by eye: the task with the most exact repeats under the standard harness for which a matching rollout exists under the comparison harness. Actions are shown as the parser received them.}
\label{fig:qualitative-ban}
\end{figure*}


\paragraph{The abstraction can discard information the model needed.}
The \emph{abstract} harness keeps the diagnosis and drops the call. For error families whose repair is fully determined by the diagnosis, such as
an unknown tool or a missing required argument, that is lossless in the sense
that matters. For
families where the model needs to see \emph{what it wrote} in order to see what
was wrong with it, the abstraction can remove the evidence along with the
hazard. We expect this to bite hardest where the failure is about a specific
value rather than a structural mistake, and Section~\ref{sec:analysis}'s
per-error-family breakdown is where to look for it.

\paragraph{The effect weakens with scale, and the interventions with it.}
The inversion shrinks monotonically with model size within every family we
tested. It follows that the value of these interventions is largest exactly
where agents are cheapest to run and smallest where they are most capable. We do
not know where, or whether, the gain crosses zero; our ladder stops at
\LargestModel\ because that is what fits in memory in float32 on the hardware
available and finishes in the time available, and we
decline to extrapolate a crossing point from a fit whose zero lies outside the
range we measured. A reader with GPUs should treat ``does this vanish at 8B?''
as the first follow-up question, and we have made it cheap to answer: the probe
needs a few hundred forward passes per model, not rollouts.

\paragraph{The dominant failure at this scale is not repetition at all.}
It is worth being clear about what the agent gets wrong most often, because it
bounds how much any harness change can buy. Across rollouts, \ProseShare\ of actions are prose where a call was expected, for instance the
model announcing that it has cancelled the meeting instead of calling
\texttt{cancel\_event}. That is a
capability limit of a 0.5B model, and no arrangement of the transcript fixes it.
A further \TruncShare\ of actions are calls cut off by our \AgentTokenCap-token generation
budget, which is our artefact rather than the model's; it applies identically in
every condition and is small enough not to move the comparison, but it does
depress absolute success on the templates with long arguments.

The interventions we study act on repetition, and repetition is a substantial
but not majority share of what goes wrong here. That is the honest ceiling on
the effect sizes in Section~\ref{sec:agent}.

\paragraph{A regime where the standard harness is fine.}
When the failed action is one the model was unlikely to produce anyway, for
instance a call it emitted only because of a formatting accident, the copying
lift starts from a very low base and the absolute repeat probability stays
small. The inversion is a statement about a ratio; it is most dangerous when the
failed action was already plausible, which is precisely the case for the
near-miss errors that dominate real agent transcripts.
